\documentclass{beamer}
\usepackage{listings}
\usepackage{minted}
\usepackage{hyperref}
\usepackage{blindtext}
\usepackage{graphicx}
\usepackage{tikz-cd}
\usepackage{animate}
\usepackage{xparse}

\newcommand\dsuml{\sum\limits}

\title{BRACU CP Workshop}
\subtitle{Day 5}
\author{
    Ahnaf Shahriar Asif\\
    Md Mubasshir Chowdhury
}

\institute{BRAC University, Dhaka}

\date{30 April 2025}

\begin{document}

\frame{\titlepage}
% Add the Table of Contents on the second slide
\begin{frame}
  \frametitle{Table of Contents}
  \tableofcontents
\end{frame}

\section{Introduction}
\begin{frame}{What will we learn ?}
\begin{itemize}
    \item<1-> Divisibility  
    \item<2-> Primes
    \item<3-> Modular Arithmetic
    \item<4-> Primality testing
    \item<5-> Number theoretic functions
\end{itemize}
\end{frame}

\section{Divisibility}
\begin{frame}{Divisibility}
\begin{itemize}
  \item<1-> Let \(a\), \(b\), and \(c\) be integers. Then:
  \begin{enumerate}[(i)]
    \item If \(b \mid a\) and \(c \mid a\), then \(bc \mid a\).
    \item If \(b \mid a\) and \(c \mid b\), then \(c \mid a\).
    \item If \(b \mid a\) and \(c \mid b\), then \(bc \mid ab\).
    \item If \(c \mid a\) and \(c \mid b\), then \(c \mid (a + b)\) and \(c \mid (a - b)\).
  \end{enumerate}
\end{itemize}
\end{frame}

\section{Finding divisors of n}
\begin{frame}{Finding divisors of n}
\begin{itemize}
  \item<1-> We can loop through all numbers and see if they divide \(n\). 
  \item<2-> Looping till $\sqrt{n}$ is enough.
  \item<3-> We can find divisors of all numbers upto $n$ in $\mathcal{O}(n \ln(n))$ time.
\end{itemize}
\end{frame}

\begin{frame}[fragile]{Code}
  \frametitle{Finding divisors of n}
  \begin{itemize}
    \begin{minted}{cpp}
vector<vector<int>> divisors(n + 1);
for (int i = 1; i <= n; i++) {
    for (int j = i; j <= n; j += i) {
        divisors[j].push_back(i);
    }
}
\end{minted}

  \end{itemize}
\end{frame}

\section{Primes}
\begin{frame}[fragile]{Primes}
  \begin{itemize}
    \item<1-> A prime $p$ is greater than 1 and only divisible by 1 and $p$.
    \item<2-> A composite number is a positive integer that is not prime.
    \item<3-> Prime numbers help us calculate a lot of things easily.
  \end{itemize}
\end{frame}

\begin{frame}[fragile]{Find number of divisors of $n$}
  \begin{itemize}
    \item<1-> Given an integer $1 \leq n \leq 10^9$, find the number of divisors of $n$.
  \end{itemize}
\end{frame}
\begin{frame}[fragile]{Solution}
  \frametitle{Solution: Number of divisors}
  \begin{itemize}
    \item<1-> We can find the prime factorization of $n$.
    \item<2-> If $n = p_1^{e_1} \cdot p_2^{e_2} \cdots p_k^{e_k}$, then the number of divisors of $n$ is given by:
    \[
      d(n) = (e_1 + 1)(e_2 + 1) \cdots (e_k + 1)
    \]
  \end{itemize}
\end{frame}


\begin{frame}[fragile]{Prime Factorization}
  \frametitle{Prime Factorization}
  \begin{itemize}
    \item<1-> We can find the prime factorization of $n$ in $\mathcal{O}(\sqrt{n})$ time.
    \item<2-> We can use a \textbf{sieve} to find all primes up to $10^6$ in $\mathcal{O}(n \log(\log(n)))$ time.
  \end{itemize}
\end{frame}

\begin{frame}[fragile]{Code}
  \frametitle{Prime Factorization code}
  \begin{minted}{cpp}
  // vector<int> primes; (contains all primes upto sqrt(n))
  vector<pair<int,int>> ppf(int n){
    vector<pair<int,int>> factors;
    for (int i = 0; primes[i] * primes[i] <= n; i++){
      int cnt = 0;
      while (n % primes[i] == 0){
        n /= primes[i];
        cnt++;
      }
      if (cnt > 0) factors.push_back({primes[i], cnt});
    }
    if (n > 1) factors.push_back({n, 1});
    return factors;
  }
  // if n = 12, it will return:
  // {{2, 2}, {3, 1}} which means 12 = 2^2 * 3^1 
  \end{minted}
\end{frame}

\begin{frame}[fragile]{Sieve}
  \frametitle{Sieve}
  \begin{itemize}
    \item<1-> We can use a sieve to find all primes up to $10^6$ in $\mathcal{O}(n \log(\log(n)))$ time.
    \item<2-> We can use the sieve to find the smallest prime factor of each number.
  \end{itemize}
\end{frame}

\begin{frame}[fragile]{Sieve Simulation}
  \frametitle{Sieve Simulation}
  
\begin{center}
\begin{tabular}{|*{10}{c|}}
\hline
1 & 2 & 3 & 4 & 5 & 6 & 7 & 8 & 9 & 10 \\
\hline
11 & 12 & 13 & 14 & 15 & 16 & 17 & 18 & 19 & 20 \\
\hline
21 & 22 & 23 & 24 & 25 & 26 & 27 & 28 & 29 & 30 \\
\hline
31 & 32 & 33 & 34 & 35 & 36 & 37 & 38 & 39 & 40 \\
\hline
\end{tabular}
\end{center}
\end{frame}

\begin{frame}[fragile]{Code}
  \frametitle{Sieve Code}
  \begin{minted}{cpp}
  bool is_composite[N];
  vector<int> primes;

  void sieve(int n){
    is_composite[1] = true;
    for(int i = 2; i <= n;i++){
      if(!is_composite[i]){
        primes.push_back(i);
        for(int j = i * i; j <= n;j+=i){
          is_composite[j] = true;
        }
      }
    }
  }
\end{minted}
\end{frame}



\section{Problem solving}
\begin{frame}{Boxes Packing}
\framesubtitle{Codeforces 903C}
\begin{itemize}
    \item<1-> You have $n$ cube-shaped boxes, each with side length $a_i$
    \item<2-> You can put box $i$ into box $j$ if $a_i < a_j$, box $i$ is not already inside another box, and box $j$ is empty
    \item<3-> Boxes can be nested multiple times
    \item<4-> A box is \emph{visible} if it is not inside any other box. Your goal is to minimize the number of visible boxes
    \item<5-> \textbf{Input:} one integer $n$ ($1 \leq n \leq 5000$), followed by a list of $n$ integers $a_1$ to $a_n$ ($1 \leq a_i \leq 10^9$)
    \item<6-> \textbf{Output:} the minimum number of visible boxes
\end{itemize}
\end{frame}

\begin{frame}{Min = GCD}
	\framesubtitle{Codeforces 2084B}
	\begin{itemize}
		\item<1-> Given an array $a$ of length $n$.
		\item<2-> You can rearrange $a$ in any order.
		\item<3-> Determine if there is an arrangement of $a$ such that there exists an integer $i (1 \leq i \leq n)$ such that
		      \begin{align*}
			      \min([a_1, a_2, \ldots, a_i]) = \gcd([a_{i+1}, a_{i+2}, \ldots, a_n])
		      \end{align*}
	\end{itemize}
\end{frame}

\begin{frame}{Min = GCD}
	\framesubtitle{Codeforces 2084B}
	\begin{block}{Sample Input}
		3 \\
		2 2 3
	\end{block}
	\begin{block}{Sample Output}
		YES
	\end{block}
\end{frame}

\begin{frame}{Min = GCD}
	\framesubtitle{Codeforces 2084B}
	\begin{block}{Sample Input}
		3 \\
		2 4 3
	\end{block}
	\begin{block}{Sample Output}
		NO
	\end{block}
\end{frame}

\begin{frame}{Min = GCD}
	\framesubtitle{Codeforces 2084B}
	\begin{block}{Sample Input}
		5 \\
		3 4 6 9 6
	\end{block}
	\begin{block}{Sample Output}
		YES
	\end{block}
\end{frame}

\begin{frame}{Min = GCD}
	\framesubtitle{Codeforces 2084B}
	\textbf{Observation:} \\
	\begin{itemize}
		\item <1-> We need to divide the array into two parts, and determine which number goes where.
		\item <2-> $\gcd{(a)} \leq \min{(a)}$.
		\item <3-> The minimum number must be in the first part lets call it $mn$
		\item <4-> The second part must consists only multiples of $mn$.
		\item <5-> The second part should contain all multiples of $mn$.
		\item <6-> Hence, the answer is YES if gcd all multiples of $mn$ (except $mn$ itself) equals to $mn$. Otherwise NO.
	\end{itemize}
\end{frame}



\NewDocumentCommand{\ceil}{s O{} m}{%
	\IfBooleanTF{#1} % starred
	{\left\lceil#3\right\rceil} % \ceil*[..]{..}
	{#2\lceil#3#2\rceil} % \ceil[..]{..}
}

\NewDocumentCommand{\floor}{s O{} m}{%
	\IfBooleanTF{#1} % starred
	{\left\lfloor#3\right\rfloor} % \floor*[..]{..}
	{#2\lfloor#3#2\rfloor} % \floor[..]{..}
}

\begin{frame}{Simple Permutation}
	\framesubtitle{Codeforces 2089A}
	\begin{itemize}
		\item<1-> Given an integer $n$. Construct a permutation $p_1, p_2, \ldots, p_n$ that follows the following rules:
		\item<2-> Define $c_i = \ceil{\frac{p_1 + p_2 + \ldots + p_i}{i}}$.
		\item<3-> Then among $c_1, c_2, \ldots, c_n$, there are must be atleast $\floor{\frac{n}{3}} - 1$ prime numbers.
	\end{itemize}
\end{frame}

\begin{frame}{Simple Permutation}
	\framesubtitle{Codeforces 2089A}
	\begin{block}{Sample Input}
		5 \\
	\end{block}
	\begin{block}{Sample Output}
		2 1 3 4 5
	\end{block}
\end{frame}

\begin{frame}{Simple Permutation}
	\framesubtitle{Codeforces 2089A}
	\begin{block}{Sample Input}
		6 \\
	\end{block}
	\begin{block}{Sample Output}
		3 2 1 4 5 6
	\end{block}
\end{frame}

\begin{frame}{Simple Permutation}
	\framesubtitle{Codeforces 2089A}
	\textbf{Hint - Bertrand's Postulate:} \\
	For each positive integer $x \geq 1$, there exists at least one prime $p$ such that $x < p < 2x$.
\end{frame}

\begin{frame}{Simple Permutation}
	\framesubtitle{Codeforces 2089A}
	\textbf{Solution:} \\
	\begin{itemize}
		\item <1-> Choose a prime number $p$ between $\floor{\frac{n}{3}}$ and $\ceil{\frac{2n}{3}}$.
		\item <2-> Construct the permutation as follows:
		      \begin{align*}
			      p, p - 1, p + 1, p - 2, p + 2, \ldots \\
		      \end{align*}
		      \text{ then fill the remaining numbers in any order.}
		\item <3-> In this way, we have $c_1 = c_3 = c_5 = \ldots = c_{2 \floor{\frac{n}{3}} - 1} = p$
	\end{itemize}
\end{frame}


\begin{frame}{Problem about GCD}
	\framesubtitle{Codeforces 2043D}
	\begin{itemize}
		\item <1-> Given three integers $l$, $r$, and $G$.
		\item <2-> Find two integers $A$ and $B$ such that:
		      \begin{align*}
			       & l \leq A \leq B \leq r         \\
			       & \gcd(A, B) = G                 \\
			       & | A - B| \text{ is maximized.}
		      \end{align*}
		      If there are multiple such pairs, output the one with the smallest $A$. \\
		      If there is no such pair, output -1 -1.
	\end{itemize}
\end{frame}

\begin{frame}{Problem about GCD}
	\framesubtitle{Codeforces 2043D}
	\begin{block}{Sample Input}
		4 12 2
	\end{block}
	\begin{block}{Sample Output}
		4 10
	\end{block}
\end{frame}

\begin{frame}{Problem about GCD}
	\framesubtitle{Codeforces 2043D}
	\begin{block}{Sample Input}
		4 8 3
	\end{block}
	\begin{block}{Sample Output}
		-1 -1
	\end{block}
\end{frame}

\begin{frame}{Problem about GCD}
	\framesubtitle{Codeforces 2043D}
	\textbf{Try solving if $G = 1$} \\
	\begin{itemize}
		\item <1-> We will try $A = l$ and $B = r$. If $\gcd(A, B) = 1$, we found our answer.
		\item <2-> Then we will try $A = l$ and $B = r - 1$, and then $A = l + 1$ and $B = r$.
		\item <3-> Then, we will try $(l, r - 2), (l + 1, r - 1), (l + 2, r)$ and so on.
		\item <4-> We will keep trying until we find a pair $(A, B)$ such that $\gcd(A, B) = 1$.
		\item <5-> If $A > B$, we will stop (There is no solution).
	\end{itemize}
\end{frame}

\begin{frame}{Problem about GCD}
	\framesubtitle{Codeforces 2043D}
	\begin{center}
		\textbf{But that seems slow?} \\
	\end{center}
\end{frame}


\begin{frame}{Problem about GCD}
	\framesubtitle{Codeforces 2043D}
	\begin{theorem}<1->
		The algorithm will find atleast one pair of coprime pair between $[l, l + 30)$ and $(r - 30, r]$.
	\end{theorem}

	\begin{proof}<2->
	\end{proof}
\end{frame}



\begin{frame}{}
\framesubtitle{}
\begin{itemize}
    \item<1-> \begin{center}
        \large{Good Bye!}
    \end{center}
\end{itemize}

\end{frame}

\end{document}





